\chapter{Case Study: Speedy Blupi and Planet Blupi}
\index{Speedy Blupi}\index{Planet Blupi}\index{case study}
\label{ch:blupi-case-study}

Two legacy Windows games shaped \texttt{free-direct}'s scope: \emph{Speedy Blupi} (through the
sibling \texttt{free-eggbert} repository) and \emph{Planet Blupi}. Both run through the
compatibility layer. A separate native-CNA feasibility study remains a proposal. The distinction
lets this chapter compare an implemented migration bridge with a source-audited port plan.

\section{Call-site-driven scope}

\texttt{free-direct}'s subsystem boundaries come from audits of the two games at named commits.
Both use DirectDraw and DirectSound; only Free Eggbert uses DirectPlay. \emph{free-eggbert} calls
\texttt{Blt} four times and \texttt{BltFast} six times; \emph{Planet Blupi} calls them three and
five times. In each game, one \texttt{Blt} performs the once-per-frame back-buffer-to-primary
present. This call-site pressure explains why Chapter~\ref{ch:freedirect-deep-dive} describes
DirectDraw as the most developed of the library's three subsystems.

\subsection{Raw occurrences versus interface calls}

A later recount reports eighteen \texttt{BltFast} occurrences in \emph{free-eggbert} and
seventeen in \emph{Planet Blupi}. That population includes header declarations, definitions of
each game's \cnaclass{CPixmap::BltFast} wrapper, and calls to both the wrapper and the DirectDraw
interface. The six and five counts include only calls that reach
\texttt{IDirectDrawSurface::BltFast}, the boundary implemented by \texttt{free-direct}. The other
eight calls per game enter the local wrapper, which then makes one of those interface calls.
Counting both levels would count the same traffic twice. Call-site totals therefore need an
explicit population definition.

\section{A pixel-format mismatch that never became a visible bug}
\label{sec:blupi-pixel-format-mismatch}

Both games request an explicit display depth: \emph{free-eggbert} selects 8 or 16 bits through a
runtime flag, while \emph{Planet Blupi} requests 8 bits. \texttt{free-direct}'s
\texttt{SetDisplayMode} ignores that value and creates a 32-bit primary surface. The deviation
does not affect these two games because their offscreen surfaces follow a separate, hard-coded
32-bit path. Every observed producer and consumer therefore agrees on the delivered format.
This is a verified harmless deviation for the audited callers, not a general guarantee about
ignoring pixel-format requests.

\section{Color keying: implementation broader than observed use}

\emph{Planet Blupi} uses source color-keying at two call sites through one caching helper. Both
construct a single-value range and match white through a \texttt{GetDC} / \texttt{SetPixel} /
\texttt{Lock} round trip, not through a fixed palette index. \texttt{free-direct} implements
low-to-high ranges, but the audited games exercise only the single-value case. A third game would
need separate evidence for ranged behavior.

\section{GetDC as a gameplay dependency}

\emph{Planet Blupi}'s \texttt{IsIconPixel} performs a per-click gameplay hit test through
\texttt{GetDC} / \texttt{ReleaseDC}. The implementation must handle both a zero-copy 32-bit wrap
and an 8-bit surface expanded into a temporary palette-converted buffer. An API that appears to
be incidental GDI interoperation therefore sits on an active gameplay path.

\section{A performance audit adds exact costs to the same call sites}

A later performance audit benchmarked the compiled library headlessly with
\texttt{SDL\_VIDEODRIVER=dummy}. DirectDraw~3's \texttt{BltFast} has no destination-size
parameter, so it cannot scale. Every audited \texttt{BltFast} interface call in both games
(\emph{free-eggbert}: \texttt{pixmap.cpp:406,574,580,612,1762,1818}; \emph{Planet Blupi}:
\texttt{pixmap.cpp:395,401,433,1235,1291}) is consequently 1:1. At the time of measurement, the
shared blit engine still performed a per-pixel integer division and bit-depth branch for these
copies. It measured
\textbf{29 times slower than a raw \texttt{memcpy}} of the identical byte count at an optimized
build (0.715~ms versus 0.0244~ms for one 640$\times$480 32-bit blit), and \textbf{206 times slower}
under the then-default unoptimized configuration.

\begin{center}
\small
\begin{tabular}{lrrr}
\toprule
\textbf{Build} & \textbf{\texttt{BltFast}} & \textbf{\texttt{memcpy}} & \textbf{Ratio} \\
\midrule
Default (no \texttt{-O} flags) & 4.704~ms & 0.0229~ms & 205.8$\times$ \\
\texttt{-DCMAKE\_BUILD\_TYPE=Release} & 0.715~ms & 0.0244~ms & 29.4$\times$ \\
\bottomrule
\end{tabular}
\end{center}

Those numbers are discovery evidence, not the current implementation. The pin now defaults an
otherwise-unspecified \emph{top-level, single-configuration} free-direct build to Release while
leaving a consuming project's choice untouched. It also detects equal source/destination extents
at runtime and performs per-row \texttt{memcpy} when format and color-key conditions permit,
followed by the required opaque-alpha fixup. The recorded follow-up benchmark reduced the same
Release case from 0.715~ms to 0.466~ms, or from 29.4 to 14.8 times the raw-copy reference. The
remaining gap is mostly the alpha pass; the original 205.8-times row no longer describes a plain
standalone configure. The two measurements preserve the before-and-after evidence instead of
presenting the discovery result as current performance.

One \texttt{Blt}-family call does require scaling: \emph{free-eggbert}'s
\texttt{CPixmap::Display()} (\texttt{pixmap.cpp:1509--1516}) presents the fixed internal
resolution into the window client area. The 1:1 fast path therefore compares the two rectangles
at runtime. The audit also found dead code on the game side: \emph{free-eggbert}'s
\texttt{CPixmap::DrawMap(int, RECT, RECT)} (\texttt{pixmap.cpp:640}) has exactly one
\texttt{Blt()} call, but no callers of its three-argument signature in the inspected source.

\section{Interfaces never exercised at all}

Across both games, the only \texttt{QueryInterface} call belongs to a DirectPlay object. None
targets the four DirectDraw classes, whose \texttt{free-direct} implementations return
\texttt{DDERR\_UNSUPPORTED}. \emph{Planet Blupi} contains no DirectPlay use, so the DirectPlay
scope described in Chapter~\ref{ch:freedirect-deep-dive} comes entirely from
\emph{free-eggbert}.

\section{DirectSound structure-layout hazard}

Both games populate \texttt{PCMWAVEFORMAT}, not the larger \texttt{WAVEFORMATEX}. The layouts
place \texttt{wBitsPerSample} at different offsets; interpreting the caller's pointer as
\texttt{WAVEFORMATEX} would read padding as the sample depth. \texttt{free-direct} therefore
stores a raw \texttt{void*} and interprets these audited calls as \texttt{PCMWAVEFORMAT}.
Both games also call \texttt{SetPan()} only on mono effects. That explains why the intended
contract was documented as mono-only, but it does not
make current panning audible. The implementation computes left/right gains and discards them;
only the ordinary volume gain reaches SDL. The tests prove clamping, return values, and no-crash
behavior, not an audible stereo split.

\section{Dead and unexercised paths}

\emph{free-eggbert}'s alternate BASS-backed sound implementation is behind a compile-time flag
fixed to false with no build override. Each game also contains a wave-loading source file that
only includes itself. Finally, every inspected call to the sound-playback method passes zero for
the looping flags. This caller-side result agrees with Chapter~\ref{ch:freedirect-deep-dive}'s
implementation-side conclusion that the two games cannot reach DirectSound looping.

Two narrower searches bounded possible additions. Excluding vendored DirectX~3 SDK headers,
neither game calls \texttt{GetCurrentPosition}; \texttt{free-direct} declares only its
\texttt{SetCurrentPosition} counterpart. The method remains unimplemented until a target caller
requires it. Neither game explicitly requests an 8-bit \texttt{CreateSurface} format either.
Together with the 32-bit primary path in \S\ref{sec:blupi-pixel-format-mismatch}, this means the
audited games do not exercise \texttt{BlitFrom}'s silent mixed-depth fallthrough. That behavior
remains a limitation, even though adding a new error route was outside the target-driven scope.

\section{Native-CNA port feasibility}

The preceding evidence concerns the \texttt{free-direct} compatibility shell. A later
\texttt{cna.md} in the inspected \emph{free-eggbert} working tree, dated 2026-07-16, studies a
different question: what would it take to port \emph{Free Eggbert} directly onto CNA's
\cnaclass{Game} / \cnaclass{GraphicsDevice} / \cnaclass{SpriteBatch} surface, replacing
\texttt{free-direct} (and \texttt{free-api} beneath it) rather than merely running through it?
That 30{,}145-byte document was untracked, so it is supplemental evidence outside Appendix~
\ref{ch:repo-map}'s immutable Git pins; the exact inspected copy has SHA-256
\texttt{87b1e3de041cdbe2af85c8c4b4faf6e74f0e518f9ff72f851f339e5419ceadb3}.
It reframes the compatibility layer as a migration
\emph{bridge}, not a migration \emph{destination} --- a distinction the analysis itself states
plainly: ``a CNA application using the \texttt{FREEDIRECT} renderer has adopted the XNA-facing API, but
it still transitively depends on \texttt{free-direct}.''

\subsection{Scope from a direct line count}

\emph{Free Eggbert} is already C++; the port does not require a language translation. The
analysis measures roughly 31{,}800 lines of C++ source and headers (a later recount in the same
document gives 31{,}773 across
\texttt{src/*.cpp} and \texttt{include/*.hpp}), with more than 21{,}000 of those lines --- the
majority --- in simulation, world, object, and event logic. It therefore recommends rewriting the
platform-services layer (graphics, input, audio, video, files, networking) while leaving the
existing simulation/data layer (\texttt{decblupi}, \texttt{decmove}, \texttt{decblock},
\texttt{decdesign}, \texttt{decio}, and the game's tables) recognizable. Preserving that layer
also preserves a direct comparison point for old and new behavior.

\subsection{Preserving the CPixmap boundary}

The analysis recommends retaining \cnaclass{CPixmap}'s public shape across its 98 direct call
sites and replacing only its implementation, with
\cnaclass{Texture2D} / \cnaclass{RenderTarget2D} / \cnaclass{SpriteBatch} objects standing in for
DirectDraw surfaces. This separates rendering equivalence from the risk of editing 98 call
sites. \texttt{SetTransparent2}'s color-key-\emph{range} semantics are singled
out as the one piece of that surface CNA's current \texttt{.cnj} metadata cannot express
directly (it represents a single exact RGB value, not an inclusive range), so a
\emph{Free Eggbert}-specific loader would need to read pixels and zero alpha across the
requested range by hand via \texttt{GetData} / \texttt{SetData} --- described as
``straightforward'' precisely because the asset sheets involved are small, not because the
technique is trivial in general.

\subsection{Preserving the 50ms simulation cadence}

\emph{Free Eggbert} updates from a 50ms multimedia timer, a fixed 20Hz simulation cadence. CNA's
\cnaclass{Game} defaults to roughly 16.67ms (60Hz). Without an explicit
\texttt{TargetElapsedTime}, the port would advance the simulation about three times too fast; see
Chapter~\ref{ch:game-loop} for the timing contract. The existing F5--F8 speed controls add another
constraint. The analysis recommends
representing them as multiple fixed 50ms simulation steps per rendered frame (2/4/8 steps)
instead of shortening \texttt{TargetElapsedTime}, preserving fixed-size simulation steps.

\subsection{MIDI integration gap}

The ten inspected music files are Standard MIDI data, and CNA's own \cnaclass{Song} loader does
not register \texttt{.mid}. The feasibility analysis correctly identifies that public integration
gap, but its explanation predates a fuller decoder audit. CNA disables FluidSynth, yet its vendored
SDL3\_mixer build leaves the Timidity MIDI decoder enabled; the lower audio layer therefore does
have a synthesis route, subject to a usable Timidity instrument configuration. What is missing is
a supported \cnaclass{Song}/Content extension mapping, asset/instrument policy, and retained test
showing that this route works in CNA rather than merely compiling below it.

The analysis gives three options: pre-render every
track to OGG/FLAC against a chosen SoundFont ahead of time (lowest implementation risk, but a
licensing/distribution question for the rendered assets and the SoundFont itself); expose and
verify a reusable CNA-level MIDI route over the available decoder (cleaner for future games, but
scope beyond what one port should decide unilaterally); or keep a
\emph{Free Eggbert}-specific MIDI service bridging through \texttt{free-api}'s own existing
TinySoundFont path while every other subsystem moves to CNA. Its staged recommendation uses the
game-specific bridge for an
initial, behavior-preserving port, option one (pre-rendered assets) for a final low-maintenance
release if asset policy permits it, and option two (a real CNA-level MIDI path) only if neither
of the other two is acceptable.

\subsection{Networking reconstruction defect}

\texttt{CNetwork::Receive} reads into a local 500-byte stack buffer but never copies the data to
the caller's destination. This source-proven defect makes the current decompiled reconstruction
an unreliable behavioral specification for networking. Future multiplayer work needs a
versioned wire format and two-process tests against intended behavior, not an attempt to preserve
this function's omission.

\subsection{Renderer progression and licensing boundary}

The recommended sequence starts on \texttt{SDL\_RENDERER} for focused portable 2D validation,
then builds the same port against \texttt{FREEDIRECT} as a differential bridge between the old
and new public APIs. An independent \texttt{OPENGLES3} run supplies evidence beyond one renderer.
Part~\ref{part:renderers} describes those identities and their implementation routes. The
feasibility document predates the renderer naming migration and calls this identity
\texttt{EASYGL}; it names the same shared implementation family, not a still-selectable public
value. \texttt{FREEDIRECT} is a migration tool, not the target architecture: a build linking it
has adopted CNA's public API but retains the DirectDraw-era dependency.

\emph{Free Eggbert} declares GPL-3.0, while CNA declares Ms-PL. Combining and distributing them
in one linked program requires a licensing review. Neither the feasibility document nor this
book makes that legal determination.

\section{Porting implications}

The call-site method produces a narrower compatibility layer than a speculative full-surface
implementation, but one aligned with the two target games. It exposed details that an API-only
reading could miss, including \texttt{GetDC}'s gameplay role and the
\texttt{PCMWAVEFORMAT} layout requirement. For a native CNA port, the same method identifies
stable seams, timing assumptions, evidence gaps, and licensing questions before code is moved.
